% Copyright (c) 2026 Eugene Joseph M. Ragasa.
% SPDX-License-Identifier: Apache-2.0
% Author: Eugene Joseph M. Ragasa
% Formal proof completed: 2026-10-19
\documentclass[11pt]{article}
\usepackage[T1]{fontenc}
\usepackage{amsmath,amssymb,booktabs,tabularx}
\usepackage[margin=1in]{geometry}
\title{Proof Dossier: Dominance Is Transitive}
\author{Eugene Joseph M. Ragasa}
\date{Formal proof completed: 2026-10-19}
\begin{document}\maketitle
\section{Identity and status}
\begin{tabularx}{\textwidth}{@{}lX@{}}\toprule
Statement ID & \texttt{BG-DECISION-THM-004} \\
Status & proved; no admissions \\
Lean declaration & \texttt{ProjectKoios.BandGauge.Decision.dominates\_trans} \\
Owning module & \texttt{ProjectKoios/BandGauge/Decision/FinitePareto.lean} \\
Axioms & none \\
\bottomrule\end{tabularx}
\section{Statement}
If $a$ dominates $b$ and $b$ dominates $c$, then $a$ dominates $c$.
\section{Argument}
Coordinatewise weak improvement composes by transitivity. Choose an objective where
$a$ is strictly better than $b$. At that objective, $b$ is weakly no worse than $c$,
so $f(a,j)<f(c,j)$. Thus the composite comparison is strict somewhere.
\section{Boundary and provenance}
This is finite-decision order infrastructure and carries no convexity or continuum-front
claim. Lean checked the exact statement; the dossier and proof were AI-assisted.
\end{document}
